\documentclass[reqno]{amsart}
\usepackage[T1]{fontenc}
\usepackage{amsmath,amssymb,amsthm}
\usepackage{xr-hyper}
\usepackage[hidelinks]{hyperref}
\externaldocument[fm-]{20261125T009500--fock-model-and-theta-intertwiners__workbook}
\externaldocument[ob-]{20261126T040100--spaces-forms-and-twists__workbook}
\externaldocument[ss-]{20261126T030400--the-seesaw-identity__workbook}
\externaldocument[s2-]{20261126T051300--two-lifts-at-a-time__workbook}
\externaldocument[s3-]{20261126T051400--lattices-cuts-and-the-ramified-forms__workbook}
\externaldocument[dl-]{20261126T030500--divisibility-of-the-composite-lift__workbook}

% Chapter number; \ref{chap} gives it.
\newcounter{chap}
% Chapter numbers, generated from 20261125T012000--grh-workbook__guide.org
% by scripts/make-chapters.py.  Do not edit: change the order of the
% index and rerun the script (./preview-site.sh and CI do this).
%
% Usage in a sheet, after \newcounter{chap}:
%   % Chapter numbers, generated from 20261125T012000--grh-workbook__guide.org
% by scripts/make-chapters.py.  Do not edit: change the order of the
% index and rerun the script (./preview-site.sh and CI do this).
%
% Usage in a sheet, after \newcounter{chap}:
%   % Chapter numbers, generated from 20261125T012000--grh-workbook__guide.org
% by scripts/make-chapters.py.  Do not edit: change the order of the
% index and rerun the script (./preview-site.sh and CI do this).
%
% Usage in a sheet, after \newcounter{chap}:
%   \input{chapters}\chapnumber{<sheet ID>}
\makeatletter
\providecommand\chapters@set[2]{\@namedef{chapters@#1}{#2}}
\chapters@set{20261125T001000}{1}% zeros-to-ramp-sums
\chapters@set{20261125T001100}{2}% inert-primes-and-estimate-b
\chapters@set{20261125T005000}{3}% galois-fourier-and-averaging
\chapters@set{20261125T005500}{4}% columns-from-a-sphere
\chapters@set{20261125T000000}{5}% slope-method
\chapters@set{20261126T030100}{6}% a-twist-for-each-prime
\chapters@set{20261125T003000}{7}% the-chain-and-the-square-root
\chapters@set{20261125T010000}{8}% rank-lemma-and-zero-estimate
\chapters@set{20261125T010500}{9}% poincare-series-and-near-orthogonality
\chapters@set{20261124T223000}{10}% elliptic-curves-for-the-proof
\chapters@set{20261125T009500}{11}% fock-model-and-theta-intertwiners
\chapters@set{20261124T223100}{12}% complex-multiplication-and-inert-reduction
\chapters@set{20261125T008000}{13}% one-cm-determinant
\chapters@set{20261124T214500}{14}% supersingular-theta-divisibility
\chapters@set{20261124T214600}{15}% the-isotropic-plane
\chapters@set{20261125T013700}{16}% pairings-on-a-formal-group-of-height-two
\chapters@set{20261125T013500}{17}% hecke-nondegeneracy-on-the-tree
\chapters@set{20261126T030200}{18}% the-bound-is-attained
\chapters@set{20261125T009000}{19}% kronecker-limit-and-chowla-selberg
\chapters@set{20261125T002000}{20}% definite-theta-lifts
\chapters@set{20261125T015500}{21}% siegel-weil-in-the-smallest-case
\chapters@set{20261125T002100}{22}% definite-theta-lifts-and-norms
\chapters@set{20261125T016500}{23}% hecke-operators-and-theta-lifts
\chapters@set{20261125T016700}{24}% divisibility-along-the-chain
\chapters@set{20261125T004000}{25}% denominators-and-sha
\chapters@set{20261125T004100}{26}% tate-twists-and-sha
\chapters@set{20261126T020430}{27}% the-minus-part-by-gauss-sums
\chapters@set{20261125T014100}{28}% congruence-modules
\chapters@set{20261125T014500}{29}% eisenstein-congruences-and-selmer-classes
\chapters@set{20261125T014700}{30}% pseudocharacters-and-the-length-bound
\chapters@set{20261125T011000}{31}% galois-representations-black-box
\chapters@set{20261126T030300}{32}% the-data
\chapters@set{20261126T020210}{33}% bessel-poincare-series-and-their-gram-operators
\chapters@set{20261126T030400}{34}% the-seesaw-identity
\chapters@set{20261126T030500}{35}% divisibility-of-the-composite-lift
\chapters@set{20261126T030600}{36}% theorem-a-short-forms
\chapters@set{20261126T030700}{37}% the-determinant-of-a-chain
\chapters@set{20261126T030800}{38}% the-cost-of-one-lift
\chapters@set{20261126T030900}{39}% denominators-of-projections-to-lifts
\chapters@set{20261126T031000}{40}% twists-with-few-congruences
\chapters@set{20261126T031100}{41}% theorem-c-and-the-main-inequality
\providecommand\chapnumber[1]{%
  \@ifundefined{chapters@#1}%
    {\PackageError{chapters}{Sheet #1 is not listed in the index}%
      {List it in 20261125T012000--grh-workbook__guide.org\MessageBreak
       and run scripts/make-chapters.py.}}%
    {\edef\thechap{\@nameuse{chapters@#1}}}}
\makeatother
\chapnumber{<sheet ID>}
\makeatletter
\providecommand\chapters@set[2]{\@namedef{chapters@#1}{#2}}
\chapters@set{20261125T001000}{1}% zeros-to-ramp-sums
\chapters@set{20261125T001100}{2}% inert-primes-and-estimate-b
\chapters@set{20261125T005000}{3}% galois-fourier-and-averaging
\chapters@set{20261125T005500}{4}% columns-from-a-sphere
\chapters@set{20261125T000000}{5}% slope-method
\chapters@set{20261126T030100}{6}% a-twist-for-each-prime
\chapters@set{20261125T003000}{7}% the-chain-and-the-square-root
\chapters@set{20261125T010000}{8}% rank-lemma-and-zero-estimate
\chapters@set{20261125T010500}{9}% poincare-series-and-near-orthogonality
\chapters@set{20261124T223000}{10}% elliptic-curves-for-the-proof
\chapters@set{20261125T009500}{11}% fock-model-and-theta-intertwiners
\chapters@set{20261124T223100}{12}% complex-multiplication-and-inert-reduction
\chapters@set{20261125T008000}{13}% one-cm-determinant
\chapters@set{20261124T214500}{14}% supersingular-theta-divisibility
\chapters@set{20261124T214600}{15}% the-isotropic-plane
\chapters@set{20261125T013700}{16}% pairings-on-a-formal-group-of-height-two
\chapters@set{20261125T013500}{17}% hecke-nondegeneracy-on-the-tree
\chapters@set{20261126T030200}{18}% the-bound-is-attained
\chapters@set{20261125T009000}{19}% kronecker-limit-and-chowla-selberg
\chapters@set{20261125T002000}{20}% definite-theta-lifts
\chapters@set{20261125T015500}{21}% siegel-weil-in-the-smallest-case
\chapters@set{20261125T002100}{22}% definite-theta-lifts-and-norms
\chapters@set{20261125T016500}{23}% hecke-operators-and-theta-lifts
\chapters@set{20261125T016700}{24}% divisibility-along-the-chain
\chapters@set{20261125T004000}{25}% denominators-and-sha
\chapters@set{20261125T004100}{26}% tate-twists-and-sha
\chapters@set{20261126T020430}{27}% the-minus-part-by-gauss-sums
\chapters@set{20261125T014100}{28}% congruence-modules
\chapters@set{20261125T014500}{29}% eisenstein-congruences-and-selmer-classes
\chapters@set{20261125T014700}{30}% pseudocharacters-and-the-length-bound
\chapters@set{20261125T011000}{31}% galois-representations-black-box
\chapters@set{20261126T030300}{32}% the-data
\chapters@set{20261126T020210}{33}% bessel-poincare-series-and-their-gram-operators
\chapters@set{20261126T030400}{34}% the-seesaw-identity
\chapters@set{20261126T030500}{35}% divisibility-of-the-composite-lift
\chapters@set{20261126T030600}{36}% theorem-a-short-forms
\chapters@set{20261126T030700}{37}% the-determinant-of-a-chain
\chapters@set{20261126T030800}{38}% the-cost-of-one-lift
\chapters@set{20261126T030900}{39}% denominators-of-projections-to-lifts
\chapters@set{20261126T031000}{40}% twists-with-few-congruences
\chapters@set{20261126T031100}{41}% theorem-c-and-the-main-inequality
\providecommand\chapnumber[1]{%
  \@ifundefined{chapters@#1}%
    {\PackageError{chapters}{Sheet #1 is not listed in the index}%
      {List it in 20261125T012000--grh-workbook__guide.org\MessageBreak
       and run scripts/make-chapters.py.}}%
    {\edef\thechap{\@nameuse{chapters@#1}}}}
\makeatother
\chapnumber{<sheet ID>}
\makeatletter
\providecommand\chapters@set[2]{\@namedef{chapters@#1}{#2}}
\chapters@set{20261125T001000}{1}% zeros-to-ramp-sums
\chapters@set{20261125T001100}{2}% inert-primes-and-estimate-b
\chapters@set{20261125T005000}{3}% galois-fourier-and-averaging
\chapters@set{20261125T005500}{4}% columns-from-a-sphere
\chapters@set{20261125T000000}{5}% slope-method
\chapters@set{20261126T030100}{6}% a-twist-for-each-prime
\chapters@set{20261125T003000}{7}% the-chain-and-the-square-root
\chapters@set{20261125T010000}{8}% rank-lemma-and-zero-estimate
\chapters@set{20261125T010500}{9}% poincare-series-and-near-orthogonality
\chapters@set{20261124T223000}{10}% elliptic-curves-for-the-proof
\chapters@set{20261125T009500}{11}% fock-model-and-theta-intertwiners
\chapters@set{20261124T223100}{12}% complex-multiplication-and-inert-reduction
\chapters@set{20261125T008000}{13}% one-cm-determinant
\chapters@set{20261124T214500}{14}% supersingular-theta-divisibility
\chapters@set{20261124T214600}{15}% the-isotropic-plane
\chapters@set{20261125T013700}{16}% pairings-on-a-formal-group-of-height-two
\chapters@set{20261125T013500}{17}% hecke-nondegeneracy-on-the-tree
\chapters@set{20261126T030200}{18}% the-bound-is-attained
\chapters@set{20261125T009000}{19}% kronecker-limit-and-chowla-selberg
\chapters@set{20261125T002000}{20}% definite-theta-lifts
\chapters@set{20261125T015500}{21}% siegel-weil-in-the-smallest-case
\chapters@set{20261125T002100}{22}% definite-theta-lifts-and-norms
\chapters@set{20261125T016500}{23}% hecke-operators-and-theta-lifts
\chapters@set{20261125T016700}{24}% divisibility-along-the-chain
\chapters@set{20261125T004000}{25}% denominators-and-sha
\chapters@set{20261125T004100}{26}% tate-twists-and-sha
\chapters@set{20261126T020430}{27}% the-minus-part-by-gauss-sums
\chapters@set{20261125T014100}{28}% congruence-modules
\chapters@set{20261125T014500}{29}% eisenstein-congruences-and-selmer-classes
\chapters@set{20261125T014700}{30}% pseudocharacters-and-the-length-bound
\chapters@set{20261125T011000}{31}% galois-representations-black-box
\chapters@set{20261126T030300}{32}% the-data
\chapters@set{20261126T020210}{33}% bessel-poincare-series-and-their-gram-operators
\chapters@set{20261126T030400}{34}% the-seesaw-identity
\chapters@set{20261126T030500}{35}% divisibility-of-the-composite-lift
\chapters@set{20261126T030600}{36}% theorem-a-short-forms
\chapters@set{20261126T030700}{37}% the-determinant-of-a-chain
\chapters@set{20261126T030800}{38}% the-cost-of-one-lift
\chapters@set{20261126T030900}{39}% denominators-of-projections-to-lifts
\chapters@set{20261126T031000}{40}% twists-with-few-congruences
\chapters@set{20261126T031100}{41}% theorem-c-and-the-main-inequality
\providecommand\chapnumber[1]{%
  \@ifundefined{chapters@#1}%
    {\PackageError{chapters}{Sheet #1 is not listed in the index}%
      {List it in 20261125T012000--grh-workbook__guide.org\MessageBreak
       and run scripts/make-chapters.py.}}%
    {\edef\thechap{\@nameuse{chapters@#1}}}}
\makeatother
\chapnumber{20261126T051500}
\renewcommand{\thesection}{\thechap.\arabic{section}}

\theoremstyle{definition}
\newtheorem{problem}{Problem}[section]
\newtheorem{definition}{Definition}[section]
\theoremstyle{plain}
\newtheorem{theorem}{Theorem}[section]

\newcommand{\Z}{\mathbb{Z}}
\newcommand{\Q}{\mathbb{Q}}
\newcommand{\C}{\mathbb{C}}
\newcommand{\F}{\mathbb{F}}
\newcommand{\Oc}{\mathcal{O}}
\newcommand{\Tc}{\mathcal{T}}
\newcommand{\Pc}{\mathcal{P}}
\newcommand{\Hom}{\operatorname{Hom}}
\newcommand{\Aut}{\operatorname{Aut}}
\newcommand{\Nm}{\mathrm{Nm}\,}
\newcommand{\Heis}{\mathrm{Heis}}
\newcommand{\vac}{\mathrm{vac}}
\newcommand{\zt}{\mathrm{zt}}
\newcommand{\gl}{\mathrm{gl}}
\newcommand{\Gl}{\mathrm{Gl}}
\newcommand{\pl}{\mathrm{pl}}
\newcommand{\res}{\mathrm{res}}
\newcommand{\ssx}{\mathrm{ss}_\xi}
\newcommand{\Cut}{\mathrm{Cut}}
\newcommand{\mult}{\mathrm{mult}}
\newcommand{\ev}{\phi^{\mathrm{ev}}}
\newcommand{\cres}{c_{\mathrm{res}}}
\newcommand{\cgl}{c_{\mathrm{gl}}}
\newcommand{\cvac}{c_{\mathrm{vac}}}
\newcommand{\z}{{\mathfrak{z}}}
\newcommand{\tl}[1]{\texttt{#1}}

\setlength{\emergencystretch}{3em}

\title{The seesaw tensor: translated vacua, gluing and the unit}
\author{}
\date{November 26, 2026}

\begin{document}
\maketitle
\refstepcounter{chap}\label{chap}

\paragraph{Prerequisites.}
From Chapter~\ref{s3-chap}: the cuts \(\Cut_j(\xi)\), their multiplicities \(\mult_j\) and the seesaw tensor \(\ssx\) (Definitions~\ref{s3-def:cut} and~\ref{s3-def:ss}), the generator ideles \(\varpi_\xi\), \(\varpi_M\) (Definition~\ref{s3-def:u}), Problems~\ref{s3-pr:mult-dvd} and~\ref{s3-pr:cut-count}, and Theorem~\ref{s3-thm:proportional}.  From Chapter~\ref{s2-chap}: the join \(\boxtimes\), the zero term \(\zt\), the stacked map \(A\) and the evaluated pair datum \(\ev_i(\xi,y)\) (Definition~\ref{s2-def:evaluated}).  From Chapter~\ref{ss-chap}: the translations \(\mathfrak t\in\Tc\), the translated vacua \(v_j(\mathfrak t)\), the data \(\mathbf 1_j(\xi)\) of the lower chain, the levels \(K^X_j\), the restriction \(\res\), chain data and the composite lift \(B_{\vartheta,\Psi}\) (Section~\ref{ss-sec:theorem}), and Problems~\ref{ss-pr:lattice}, \ref{ss-pr:unit}, \ref{ss-pr:degenerate}.  In a tensor over pairs, later pairs stand on the left.  From Chapter~\ref{fm-chap}: theta lines and intertwiners (Problems~\ref{fm-pr:G-line}, \ref{fm-pr:I-explicit}, \ref{fm-pr:p-neighbours}, Theorem~\ref{fm-thm:compose}).  From Chapter~\ref{ob-chap}: the frames of the spaces, \(Q_\xi\), \(\Nm Q_\xi\), the integral vacuum and \(\cvac\) (Sections~\ref{ob-sec:spaces}, \ref{ob-sec:twists}, \ref{ob-sec:vacuum}), and the ring \(\Oc^{[q]}\) of the \(z\in\C\) with \(\lvert\eta(z)\rvert_\ell\le1\) for every prime \(\ell\nmid q\) and every ring map \(\eta\colon\C\to\C_\ell\).  The Weil representation over the finite adeles in its Schr\"odinger and lattice models \cite{We}, \cite[Ch.~2]{MVW}.  As in Chapter~\ref{s2-chap}, \(\mathbf 1_j(\xi)\) is the finite datum of the lift \(Y_{j-1}\to Y_j\); \(\omega^X_j\) is the finite Weil action on the data of the lift \(X_j\to X_{j+1}\).

\medskip
This chapter is the last of the four that prove the seesaw identity, Theorem~\ref{ss-thm:seesaw}, which Theorem~A uses.  Theorem~\ref{s3-thm:proportional} gives it on automorphic forms, for ramified source forms, with the tensor \(\ssx\), a weighted sum over the cuts of evaluated pair data, in place of a chain datum.  Which scalar and which common denominator suffice to express \(\ssx\) in glued vacuum tensors \(\gl(\mathfrak t)\)?  Three local facts answer this: Fourier inversion on a residue module, the lattice model of the vacuum, and a neighbour lattice at \(N\).  One complex unit then works for all twists.

Throughout, \(\xi\) is a tame character and \(q:=\Nm Q_\xi\).  A finite adelic vector is \emph{integral away from \(q\)} if its entries are integral at every finite place not above a prime dividing \(q\).

\section{Residue modules and Fourier inversion}\label{sec:residues}

\begin{definition}\label{def:combination}
Let \((w_\gamma)_{\gamma\in\mathcal Z}\) be vectors of a complex vector space, \(\mathfrak d\ge1\) an integer and \(\mathfrak s\ge0\) real.  A vector \(z\) is a \emph{combination of the \(w_\gamma\), \(\gamma\in\mathcal Z\), with denominator \(\mathfrak d\) and size \(\mathfrak s\)} if \(z=\sum_{\gamma\in\mathcal Z}r_\gamma w_\gamma\) with finitely many nonzero \(r_\gamma\in\C\), \(\mathfrak d\,r_\gamma\in\Oc^{[q]}\) for every \(\gamma\) and, if \(Q_\xi=\varnothing\), \(\sum_\gamma\lvert\sigma(r_\gamma)\rvert\le\mathfrak s\) for every \(\sigma\in\Aut\C\).
\end{definition}

Let \(\mathfrak m\subseteq\Oc_E\) be a nonzero ideal, \(\mathsf A_{\mathfrak m}:=(\Oc_E/\mathfrak m)^n\), \(\psi_{\mathfrak m}\) a primitive additive character of \(\Oc_E/\mathfrak m\), \(\mathfrak y\cdot\mathfrak b:=\sum_k\mathfrak y_k\mathfrak b_k\), and \(\chi_{\mathfrak b}(\mathfrak y):=\psi_{\mathfrak m}(\mathfrak y\cdot\mathfrak b)\).  For the next problem alone, take the additive group \(\Z/5\Z\) and characters with values in the fifth roots of unity of \(\F_{11}\).

\begin{problem}\label{pr:fourier}
The class of \(3\) in \(\F_{11}\) has order \(5\), and \(\sum_{l=0}^{4}3^{\,l(j-2)}\) is \(0,0,5,0,0\) in \(\F_{11}\) for \(j=0,\dots,4\).
\end{problem}

At the places of \(\mathfrak m:=(\pi_M)\prod_{w\in Q_\xi}w\) (a modulus prime to its conjugate) the datum \(\mathbf 1_j(\xi)\) is the operator of the Schr\"odinger model attached to the point mass at the class \(\mathfrak i\in \mathsf A_{\mathfrak m}\) of the inclusion matrix, \(n=\dim Y_{j-1}\dim Y_j\), applied to the vacuum \(\vac\); at the other finite places it is the vacuum.  Here \(\psi_{\mathfrak m}(u\bmod\mathfrak m)=\psi_0(\varpi_\xi^{-1}u)\) for integral finite adeles \(u\), with \(\varpi_\xi\) the generator idele of \(\mathfrak m\) of Definition~\ref{s3-def:u} and \(\psi_0\) the self-dual character of the finite adeles; a residue function is read on the reduction at the places of \(\mathfrak m\), and is zero off the integral box there.  The operator of \(\chi_{\mathfrak b}\), on vectors fixed by the integral translations at the places of \(\mathfrak m\) and their conjugates (as \(\vac\) is), is the Heisenberg translation \(\Heis(t_{\mathfrak b})\), where \(t_{\mathfrak b}\) is the complex conjugate of the part at the places of \(\mathfrak m\) of \(\varpi_\xi^{-1}\tilde u_{\mathfrak b}\), \(\tilde u_{\mathfrak b}\) any integral lift of \(\mathfrak b\).

\begin{problem}\label{pr:delta}
The point mass at \(\mathfrak i\) is \(\lvert \mathsf A_{\mathfrak m}\rvert^{-1}\sum_{\mathfrak b}\psi_{\mathfrak m}(-\mathfrak i\cdot\mathfrak b)\,\chi_{\mathfrak b}\).  Hence \(\mathbf 1_j(\xi)=\lvert \mathsf A_{\mathfrak m}\rvert^{-1}\sum_{\mathfrak b}\psi_{\mathfrak m}(-\mathfrak i\cdot\mathfrak b)\,\Heis(t_{\mathfrak b})\vac\), and for every \(\sigma\in\Aut\C\) the absolute values of the \(\sigma\)-conjugates of these coefficients sum to \(1\).
\end{problem}

\begin{problem}\label{pr:delta-den}
Each \(\bar\pi_Mt_{\mathfrak b}\) is integral away from \(q\), and \(\#(\Oc_E/\pi_M\Oc_E)^n\) divided by \(\lvert \mathsf A_{\mathfrak m}\rvert\) is \(q^{-n}\in\Oc^{[q]}\).  So \(\mathbf 1_j(\xi)\) is a combination of the \(\Heis(\z)\vac\), over the \(\z\) with \(\bar\pi_M\z\) integral away from \(q\), with denominator \(\#(\Oc_E/\pi_M\Oc_E)^n\) and size \(1\).
\end{problem}

\section{The lattice model}\label{sec:lattice}

Let \(V\) be \(\Hom(Y_{j-1},Y_j)\) or \(\Hom(X_j,X_{j+1})\), with its standard lattice \(\Lambda\) in the finite adelic frame of Section~\ref{ob-sec:spaces}, finite Heisenberg translations \(\Heis(\z)\), Weil representation \(\omega\) and vacuum \(\vac\).  For background on the Weil representation and its models see \cite{We}, \cite[Ch.~2]{MVW}; we use the following, in the normalization of \cite{GL}.  (L1)~\(\Heis(\z)\vac=\vac\) for \(\z\in\Lambda\), and the vectors of the finite part fixed by all these translations are the multiples of \(\vac\) (compare Problem~\ref{fm-pr:G-line}).  (L2)~\(\Heis(\z)\Heis(\z')=\zeta\,\Heis(\z+\z')\), where \(\zeta\) is a root of unity depending only on \(\sum_k\bar\z'_k\z_k\) and equal to \(1\) when this sum is \(0\).  (L3)~\(\omega(h)\Heis(\z)=\Heis(h\z)\,\omega(h)\) for \(h\in U(V)(\mathbb A_f)\).  (L4)~\(\omega(h)\vac=c_h\sum_{\z\in\Lambda/(\Lambda\cap h^{-1}\Lambda)}\Heis(h\z)\vac\) with \(c_h\ne0\); compare the intertwiner defined in Section~\ref{fm-sec:intertwiners} and evaluated in Problem~\ref{fm-pr:I-explicit}.  The scalar \(c_h\) is the Weil multiplier divided by its squared absolute value.  If \(h\) is integral away from \(q\), then \(\lvert\eta(c_h)\rvert_\ell=1\) for every prime \(\ell\nmid q\) and every ring map \(\eta\colon\C\to\C_\ell\); if \(h\) is integral at every finite place, then \(c_h=1\).  \emph{Status}: we use the classical lattice model in the normalization of \cite{GL}, which proves (L1)--(L4) and the multiplier bounds.

\begin{problem}\label{pr:line}
Let \(h\in U(V)(\mathbb A_f)\).  A vector \(f\) of the finite part with \(\Heis(h\z)f=f\) for all \(\z\in\Lambda\) is a multiple of \(\omega(h)\vac\).
\end{problem}

\begin{problem}\label{pr:move}
Let \(h\in U(V)(\mathbb A_f)\) be integral away from \(q\), and let \(\bar\pi_M\z_0\) be integral away from \(q\).  Then \(\omega(h)\Heis(\z_0)\vac\) is a combination of the \(\Heis(\z)\vac\), over the \(\z\) with \(\bar\pi_M\z\) integral away from \(q\), with denominator \(1\) and size \(1\).
\end{problem}

The action \(\omega_j\) on the data of the lower chain is that of Problem~\ref{pr:move} times a determinant-character factor, which at points integral away from \(q\) has \(\ell\)-adic norm \(1\) for \(\ell\nmid q\) and absolute value \(1\) under every \(\sigma\in\Aut\C\).  So the moved data \(\mathbf 1_j(\xi)^y\) at points integral away from \(q\) keep the denominators of Problem~\ref{pr:delta-den} and size \(1\).  If the two translations become integral away from \(q\) after multiplication by \(\bar\pi_M\), the join of their translated vacua is a root of unity times a translated vacuum on \(\Hom(Y_{2i+1},Y_{2i}^-\oplus Y_{2i+2})\), whose translation becomes integral away from \(q\) after multiplication by \(\pi_M\) on the first block and by \(\bar\pi_M\) on the second.  Evaluation translates the joined datum by \(A_f\), takes its finite zero term along the isotropic part \(I\) of the splitting, and applies \(\boxtimes^{-1}\) on the definite spaces.

\begin{problem}\label{pr:reduction}
Let \(F\) be a nonarchimedean local field with integers \(\Oc_F\), \(W\) a finite-dimensional \(F\)-vector space with a nondegenerate symmetric or alternating bilinear form, \(L\subset W\) a lattice with \(L=\{x:\langle x,L\rangle\subseteq\Oc_F\}\), and \(I\subseteq W\) an isotropic subspace.  Then the image of \(L\cap I^\perp\) in \(I^\perp/I\) is a lattice equal to its own dual.
\end{problem}

Let \(L\) be the standard lattice of \(\Hom(Y_{2i+1},Y_{2i}^-\oplus Y_{2i+2})\) in its finite adelic frame.  The proof identifies the reduced lattice with a finite adelic unitary image of the standard lattice in the quotient and deduces its self-duality for the pairing \((x,w)\mapsto\psi_0(\bar x^tw)\); Problem~\ref{pr:reduction} is the analogue at one place.  The image of \(L\cap I^\perp\) in \(I^\perp/I\), identified by the splitting with the finite adelic space of \(\Hom(X_{2i+1},X_{2i}^-\oplus X_{2i+2})\), is the \emph{zero-term lattice} of pair \(i\).  The zero term of the vacuum on the first Hom space is a number \(c_i'\ne0\) times the sum of the translates of the vacuum on the second by representatives of the zero-term lattice modulo its intersection with the standard lattice there.

\section{The neighbour at \texorpdfstring{\(N\)}{N}}\label{sec:neighbour}

Which lattice is the zero-term lattice?  Let \(P\) be the set of places \(w\mid\pi_N\) and \(\bar P\) the set of their conjugates; they are disjoint, with union the set of places above \(N\).  Let \(\mathbf 1_P\) be the finite adele equal to \(1\) on \(P\) and \(0\) elsewhere, \(\pi_P\) the finite idele equal to \(\pi_N\) on \(P\), to \(\bar\pi_N^{-1}\) on \(\bar P\) and to \(1\) elsewhere, and \(u_{-1}\) a fixed unit of the integral finite adeles with \(u_{-1}\bar u_{-1}=-1\).  For a pair \(i\) let \(e\), \(e''\) be the last basis vectors of \(X_{2i}\) and \(X_{2i+2}\) in \(X_{2i}^-\oplus X_{2i+2}\), and put \(e_P:=e+\bar u_{-1}\bar{\mathbf 1}_Pe''\).  For columns write \(z^*:=\bar z^t\).

\begin{problem}\label{pr:axis}
Over a commutative ring with involution, let \(z\) be a column with \(z^*z=1\), \(\delta\) a scalar and \(\gamma:=1+(\delta-1)zz^*\).  Then \(\gamma z=\delta z\), \(\gamma z'=z'\) when \(z^*z'=0\), and \(\gamma^*\gamma=1\) if and only if \(\delta\bar\delta=1\).  Moreover \(e_P^*e_P=1\) and \(\pi_P\bar\pi_P=1\).
\end{problem}

\begin{definition}\label{def:neighbour}
Define the \emph{neighbour lattice} of pair \(i\) as the image of the standard lattice of \(\Hom(X_{2i+1},X_{2i}^-\oplus X_{2i+2})\) under the matrix acting on the target as \(\gamma_P:=1+(\pi_P-1)e_Pe_P^*\) at the places above \(N\), and at the other finite places as the identity for \(i<\varrho\); for \(i=\varrho\), at these places it scales the last row of \(X_{m}\) by \(\varpi_M/\bar\varpi_M\) and fixes the other rows, \(\varpi_M\) the generator idele of \((\pi_M)\) of Definition~\ref{s3-def:u}.
\end{definition}

For a local example take \(N=5\) and the split algebra \(\Q_5\times\Q_5\), with factors indexed by \(w\in P\) and \(\bar w\) and conjugation exchanging them; take \(\pi_N=(5,1)\), \(u_{-1}=(1,-1)\), \(\mathbf 1_P=(1,0)\), and keep the rows \(e\), \(e''\).

\begin{problem}\label{pr:split-neighbour}
In this model \(e_P=\bigl((1,1),(0,1)\bigr)\), \(\pi_P=(5,\tfrac15)\), and the components of \(\gamma_P\) at \(w\) and \(\bar w\) are \(G=\left(\begin{smallmatrix}5&4\\0&1\end{smallmatrix}\right)\) and \(G'=\left(\begin{smallmatrix}1/5&0\\-4/5&1\end{smallmatrix}\right)\), with \(G^tG'=1\).
\end{problem}

To a residue modulo \(N\) at \(w\), with representative \(y\in\{0,\dots,N-1\}\), the proof attaches the translation supported at \(\bar w\) with entries \(u_{-1}\bar y/N\) in row \(e\) and \(-\bar y/N\) in row \(e''\), where \(\bar y/N\) is the conjugate of the adele equal to \(y/N\) at \(w\) and \(0\) elsewhere.

\begin{problem}\label{pr:one-residue}
\(G\) is integral, and \(G'\Z_5^2+\Z_5^2=\Z_5^2+\Z_5\cdot\tfrac15(1,1)^t\); the five classes modulo \(\Z_5^2\) are those of \(\tfrac15(y,y)^t\), \(y=0,\dots,4\).  The translation attached to \(y\) is \(-\tfrac15(y,y)^t\) at \(\bar w\), so one residue labels both rows.
\end{problem}

\emph{Black box.}  For every pair \(i\) of the proof the zero-term lattice is the neighbour lattice; the lift of its columns is integral and orthogonal to the finite stacked map.  We take the lattice equality and the two properties of the lift from \cite{GL}.

A representative of the neighbour lattice modulo its integral points is indexed by one residue modulo \(N\) for each basis vector of \(X_{2i+1}\) and each \(w\in P\), specifying both rows \(e\), \(e''\), and one class away from \(N\).  Under the join on the definite spaces, \(\Gl_i\) corresponds to the sum of the Heisenberg operators of these residue translations.  Let \(C_i\) be \(c_i'\) times the number of classes away from \(N\); it depends only on the data.  The proof obtains: for \(n\) prime to \(N\) and \(\z\) on \(\Hom(Y_{2i+1},Y_{2i}^-\oplus Y_{2i+2})\) with \(\pi_M\) times its first block and \(\bar\pi_M\) times its second integral away from \(n\), the evaluation of \(\Heis(\z)\vac\) is a finite sum of the vectors \(C_i\,\Gl_i(\Heis(t_1)\vac\otimes\Heis(t_2)\vac)\), with \(\bar\pi_Mt_1\), \(\bar\pi_Mt_2\) integral away from \(n\), whose coefficients become integral away from \(n\) after multiplication by \(\#(\Oc_E/\pi_M\Oc_E)^{(n_{2i}+n_{2i+2})n_{2i+1}}\) and have \(\sigma\)-conjugates of absolute sum at most \(1\) for every \(\sigma\in\Aut\C\).

\section{Invariance}\label{sec:joint}

\begin{problem}\label{pr:swap}
Let \(\mathrm{sw}\) exchange the coordinates of \(\Q^2\).  For \(v=(1,2)^t\) and \(v'=(3,-1)^t\), the sum of the matrices in the orbit of \(v{v'}^t\) under \(X\mapsto\mathrm{sw}\,X\,\mathrm{sw}^t\) is \(\left(\begin{smallmatrix}1&5\\5&1\end{smallmatrix}\right)\) and is fixed by this action.
\end{problem}

\begin{problem}\label{pr:joint}
Let a group \(H\) act on finite sets \(\mathcal X_0,\dots,\mathcal X_{m-1}\) and linearly on a complex vector space; let \(\mult_j\colon\mathcal X_j\to\Z_{>0}\) be \(H\)-invariant and \(\mathcal E\) a function on \(\mathcal X_0\times\dots\times\mathcal X_{m-1}\) with \(\mathcal E(bc)=b\,\mathcal E(c)\) for \(b\in H\).  Then \(\sum_c\prod_j\mult_j(c_j)^{-1}\mathcal E(c)\) is fixed by \(H\).
\end{problem}

In the proof, let \(b_j\in K^X_j\) for \(j<m\) and \(b_m:=1\).  For \(j<m\) the image of \(b_j\) in \(U(Y_j)(\mathbb A_f)\) lies in the finite Bessel subgroup, permutes \(\Cut_j(\xi)\) and keeps the multiplicities; and evaluation is equivariant: the evaluated pair datum at the translated points is \(\omega^X_{2i}(b_{2i},b_{2i+1})\otimes\omega^X_{2i+1}(b_{2i+1},b_{2i+2})\) applied to \(\ev_i(\xi,y)\).  By Problem~\ref{pr:joint}, \(\ssx\) is fixed by \(\bigotimes_j\omega^X_j(b_j,b_{j+1})\) for all such \(b\).

\section{Denominators and sizes}\label{sec:theorem}

\begin{definition}\label{def:cres}
Define the \emph{residue count} by \(\cres:=M^{fn_m^2}\).
\end{definition}

Since \(q\) is prime to \(N\), the pair evaluation applies with \(n=q\).  Thus for \(0\le i\le\varrho\) and points \(y\) integral away from \(q\), \(\ev_i(\xi,y)\) is a combination of the vectors \(C_i\,\Gl_i(\Heis(t_1)\vac\otimes\Heis(t_2)\vac)\), over \(t_1,t_2\) with \(\bar\pi_Mt_1\), \(\bar\pi_Mt_2\) integral away from \(q\), with denominator \(\#(\Oc_E/\pi_M\Oc_E)^{\mathfrak e_i}\) and size \(1\), where \(\mathfrak e_i:=\dim Y_{2i}\dim Y_{2i+1}+\dim Y_{2i+1}\dim Y_{2i+2}+(n_{2i}+n_{2i+2})n_{2i+1}\).  The proof states this with the uniform denominator \(\cres^4\).

\begin{problem}\label{pr:pair-exponent}
Substituting \(1,2,3\) for the three \(Y\)-dimensions and for the three \(X\)-ranks in \(\mathfrak e_i\) gives \(16\le4\cdot3^2\).  For the ranks of the proof, \(\mathfrak e_i\le4n_m^2\), so \(\#(\Oc_E/\pi_M\Oc_E)^{\mathfrak e_i}\) divides \(\cres^4\).
\end{problem}

\begin{problem}\label{pr:substitute}
Let \(z\) be a combination of the \(w_\gamma\), \(\gamma\in\mathcal Z\), with denominator \(\mathfrak d\) and size \(\mathfrak s\), and let each \(w_\gamma\) be a combination of the \(w'_{\gamma'}\), \(\gamma'\in\mathcal Z'\), with denominator \(\mathfrak d'\) and size \(\mathfrak s'\).  Then \(z\) is a combination of the \(w'_{\gamma'}\) with denominator \(\mathfrak d\mathfrak d'\) and size \(\mathfrak s\mathfrak s'\).  If \(z_k\) is a combination of the \(w_{k,\gamma}\) with denominator \(\mathfrak d_k\) and size \(\mathfrak s_k\) \((k<l)\), then \(z_{l-1}\otimes\dots\otimes z_0\) is a combination of the \(w_{l-1,\gamma_{l-1}}\otimes\dots\otimes w_{0,\gamma_0}\) with denominator \(\prod_k\mathfrak d_k\) and size \(\prod_k\mathfrak s_k\).
\end{problem}

The integral vacuum of \(\Hom(X_j,X_{j+1})\) is \(\cvac^{\,n_jn_{j+1}}\vac\) (Section~\ref{ob-sec:vacuum}).  Put
\[
  \cgl:=\prod_{i=0}^{\varrho}\frac{C_i}{\cvac^{\,n_{2i}n_{2i+1}}\,\cvac^{\,n_{2i+1}n_{2i+2}}} ,
\]
a nonzero number.  By Problems~\ref{pr:pair-exponent} and~\ref{pr:substitute}, for points \(y\) integral away from \(q\) the tensor \(\bigotimes_{i=0}^{\varrho}\ev_i(\xi,y)\) is a combination of the \(\cgl\,\gl(\mathfrak t)\), over the \(\mathfrak t\) with every \(\bar\pi_M\mathfrak t_j\) integral away from \(q\), with denominator \(\cres^{2m}\) and size \(1\).

\begin{problem}\label{pr:cut-weights}
Each \(\mult_j(c)^{-1}\), \(c\in\Cut_j(\xi)\), lies in \(\Oc^{[q]}\), and if \(Q_\xi=\varnothing\), then \(\sum_{c_0,\dots,c_{m-1}}\prod_j\mult_j(c_j)^{-1}\le\cres^m\).  Hence \(\ssx\) is a combination of the \(\bigotimes_i\ev_i(\xi,c)\), over tuples \(c\) of cut points with representatives integral away from \(q\), with denominator \(1\) and size \(\cres^m\).
\end{problem}

The pair evaluation, the tensor bound above and Problem~\ref{pr:cut-weights}, followed by substitution (Problem~\ref{pr:substitute}), give the following theorem.

\begin{theorem}\label{thm:tensor}
There is \(\cgl\in\C\), depending only on the data, such that for every tame character \(\xi\), with \(q=\Nm Q_\xi\), the tensor \(\ssx\) is a combination of the \(\cgl\,\gl(\mathfrak t)\), over the \(\mathfrak t\in\Tc\) with every \(\bar\pi_M\mathfrak t_j\) integral away from \(q\), with denominator \(\cres^{2m}\) and size \(\cres^m\): there are \(r_{\xi,\mathfrak t}\in\C\), finitely many nonzero, with
\begin{gather*}
  \ssx=\sum_{\mathfrak t}r_{\xi,\mathfrak t}\,\cgl\,\gl(\mathfrak t),\qquad
  \cres^{2m}r_{\xi,\mathfrak t}\in\Oc^{[q]},\\
  Q_\xi=\varnothing\ \Longrightarrow\ \sum_{\mathfrak t}\lvert\sigma(r_{\xi,\mathfrak t})\rvert\le\cres^m\ \text{for all }\sigma\in\Aut\C .
\end{gather*}
\end{theorem}

\section{The unit and the seesaw identity}\label{sec:unit}

By Theorem~\ref{s3-thm:proportional} and Theorem~\ref{thm:tensor}, for every ramified source form \(\varphi\) the composite of the lower lifts on automorphic forms applied to \(\Pc_{g,\lambda}\varphi\) is \(c'\cgl/\cres^{2m}\) times \(\res\,\Theta^X_\xi\bigl(\sum_{\mathfrak t}\cres^{2m}r_{\xi,\mathfrak t}\gl(\mathfrak t)\bigr)\varphi\), and these coefficients lie in \(\Oc^{[q]}\).  The scalar may be \(0\).

\begin{problem}\label{pr:unit-check}
Let \(c'=2\), \(\cgl=3\), denominator \(4\) and \(r=(\tfrac14,-\tfrac12)\).  Applying Problem~\ref{ss-pr:unit} to the scalar \(c'\cgl/4\) gives \(u=\tfrac23\) and \(c=u^{-1}=\tfrac32\).  The cleared vector \(\Psi=4r\) is \((1,-2)\) and satisfies \(c\Psi=c'\cgl r=(\tfrac32,-3)\).  If \(c'\cgl=0\), the same identity holds with \(c=1\) and \(\Psi=0\).
\end{problem}

\begin{problem}\label{pr:truncate}
Let \(\mathcal V\) be the set of finite places above the primes dividing \(Mq\), and \(\mathbf 1_{\mathcal V}\mathfrak t:=(\mathbf 1_{\mathcal V}\mathfrak t_j)_j\).  If every \(\bar\pi_M\mathfrak t_j\) is integral away from \(q\), then \(v_j(\mathbf 1_{\mathcal V}\mathfrak t)=v_j(\mathfrak t)\) for every \(j\), so \(\gl(\mathbf 1_{\mathcal V}\mathfrak t)=\gl(\mathfrak t)\).
\end{problem}

\begin{problem}\label{pr:chain-datum}
Let the tame character \(\xi\) be a twist \(\vartheta\), with \(\cgl\) as in Theorem~\ref{thm:tensor}, coefficients \(r_{\vartheta,\mathfrak t}\) for this choice in Theorem~\ref{thm:tensor}, and \(\mathcal V\) as in Problem~\ref{pr:truncate} for \(q=\Nm Q_\vartheta\).  If \(c'\cgl\ne0\), put \(\Psi(\mathfrak t):=\cres^{2m}\sum r_{\vartheta,\mathfrak t'}\) over the \(\mathfrak t'\) with \(\mathbf 1_{\mathcal V}\mathfrak t'=\mathfrak t\); otherwise \(\Psi:=0\).  Then \(\Psi\) is a chain datum for \(\vartheta\); for \(Q_\vartheta=\varnothing\) its size is at most \(\cres^{2m}\cres^{m}=M^{3fmn_m^2}\).
\end{problem}

Theorem~\ref{s3-thm:proportional}, Theorem~\ref{thm:tensor}, the invariance of \(\ssx\), Problems~\ref{ss-pr:unit}, \ref{pr:truncate} and~\ref{pr:chain-datum}, and the passage from automorphic to algebraic forms on \(Y_m\) prove Theorem~\ref{ss-thm:seesaw}, with \(c=c'\cgl/\cres^{2m}\) if \(c'\cgl\ne0\) and \(c=1\), \(\Psi=0\) otherwise (compare Problem~\ref{ss-pr:degenerate}).  Passage to algebraic forms on \(Y_m\) carries the restriction on the right of the automorphic identity to the restriction in \(B_{\vartheta,\Psi}\).  Nothing is said about the size of \(c\).

\begin{definition}\label{def:plain}
For \(\mathfrak t\in\Tc\) let \(\pl_0(\mathfrak t):=1\in\C\) and \(\pl_{i+1}(\mathfrak t):=\bigl(v_{2i}(\mathfrak t)\otimes v_{2i+1}(\mathfrak t)\bigr)\otimes\pl_i(\mathfrak t)\) for \(0\le i\le\varrho\).  Define the \emph{plain vacuum tensor} by \(\pl(\mathfrak t):=\pl_{\varrho+1}(\mathfrak t)\).
\end{definition}

If every \(\mathfrak t_j\) vanishes at the places above \(N\), then \(\gl(\mathfrak t)\) is the sum of the \(\pl(\mathfrak t+\mathfrak g_y)\) over the residues \(y\) of the gluings, \(\mathfrak g_y\) the tuple of their translations; Chapter~\ref{dl-chap} uses this expansion of the composite lift.

This proves Theorem~\ref{thm:tensor}, and with it the seesaw identity, from the statements taken as given: the Schr\"odinger and lattice models with the multiplier bounds, the zero-term lattice, its lift and the pair evaluation quoted above, the passage from automorphic to algebraic forms on \(Y_m\), and the statements marked ``in the proof''.

\paragraph{In the Lean formalization.}
{\raggedright Declaration names omit the prefix \tl{New.}.  Theorem~\ref{thm:tensor} is \tl{Data.exists\_isAdmissibleCombination\_seesawDatum} in the Lean proof \cite{GL}.  Definition~\ref{def:combination} specializes \tl{Seesaw.IsAdmissibleCombination} to \(q=\Nm Q_\xi\), the size condition \(Q_\xi=\varnothing\), positive denominators and nonnegative bounds.  Definition~\ref{def:cres} is \tl{Data.residueMatrixCount}; Definition~\ref{def:plain} and \(\gl(\mathfrak t)\) are \tl{Data.ungluedVacuum} and \tl{Data.gluedVacuum}.  The uniform bound for one evaluated pair is \tl{Data.exists\_isAdmissibleCombination\_evaluatedDatum}, the bound for their tensor \tl{Data.exists\_isAdmissibleCombination\_gluedVacuum}, and joint invariance \tl{Data.isInvariant\_seesawDatum}.  The zero-term lattice equality is \tl{Seesaw.DoubleLiftData.zeroTermLattice\_splitting\_eq}.  Theorem~\ref{ss-thm:seesaw} is \tl{Data.seesaw\_identity}.  The other statements read from the proof are, in \cite{GL} (paths relative to \path{grh/src/}): for the residue operator \path{Lib/Weil/DualPair/Defs.lean} (587, 625) and \path{New/UpperBound/Objects/Rows.lean} (377); for (L1)--(L4) \path{Lib/Weil/Construction/VacuumProjections.lean} (102, 203, 450, 755, 928), with the multiplier bounds in \path{New/UpperBound/Seesaw/AdmissibleCombinations.lean} (451) and \path{Lib/Weil/Construction/FiniteMultiplier.lean} (948, 979).  In \path{New/UpperBound/Seesaw/}: the moved and joined data \path{GoodDatum.lean} (267, 370); the evaluation, the reduced lattice and the equivariance \path{EdgePairEvaluation.lean} (487, 514, 853, 988); the zero term of the vacuum, the lift and the pair evaluation \path{AdmissibleEvaluation.lean} (248, 329, 565, 596); the residue translations and representatives \path{NeighbourLattice.lean} (171, 996); the gluing \path{NeighbourVacuum.lean} (241); the rank bound \path{GoodDatum.lean} (421); the inputs to invariance \path{Cut.lean} (173, 182, 199).  The plain expansion is in \path{New/UpperBound/Divisibility/ChainDatumExpansion.lean} (233).\par}

\section*{Hints}

\subsection*{\S\ref{sec:residues}}
\begin{description}
\item[\ref{pr:delta}] Orthogonality of the characters of \(\mathsf A_{\mathfrak m}\); the operator depends linearly on the function.  Compare Problem~\ref{ss-pr:lattice}.
\item[\ref{pr:delta-den}] \(\mathfrak m\) is the product of the coprime ideals \((\pi_M)\) and \(\prod_{w\in Q_\xi}w\).
\end{description}

\subsection*{\S\ref{sec:lattice}}
\begin{description}
\item[\ref{pr:line}] Move the problem to the standard lattice.
\item[\ref{pr:move}] (L2)--(L4).  If \(Q_\xi=\varnothing\), then \(h\Lambda=\Lambda\).
\item[\ref{pr:reduction}] Compute the annihilator of \(L\cap I^\perp\) in \(W\); \(L\cap I^\perp\) is a direct summand of \(L\).
\end{description}

\subsection*{\S\ref{sec:neighbour}}
\begin{description}
\item[\ref{pr:axis}] Expand \(\gamma^*\gamma\).  \(\mathbf 1_P\bar{\mathbf 1}_P=0\).
\end{description}

\subsection*{\S\ref{sec:joint}}
\begin{description}
\item[\ref{pr:joint}] Substitute \(c\mapsto bc\).
\end{description}

\subsection*{\S\ref{sec:theorem}}
\begin{description}
\item[\ref{pr:pair-exponent}] All ranks are at most \(n_m\), and \(\#(\Oc_E/\pi_M\Oc_E)=M^f\).
\item[\ref{pr:substitute}] \(\Oc^{[q]}\) is a ring; the triangle inequality gives the sizes.
\item[\ref{pr:cut-weights}] Problems~\ref{s3-pr:mult-dvd} and~\ref{s3-pr:cut-count}.
\end{description}

\subsection*{\S\ref{sec:unit}}
\begin{description}
\item[\ref{pr:truncate}] Split \(\mathfrak t_j\) into its parts on \(\mathcal V\) and off \(\mathcal V\).  Which part fixes the vacuum, and what is their Heisenberg phase?
\item[\ref{pr:chain-datum}] Problems~\ref{pr:truncate} and~\ref{pr:joint}; regroup equal indices.
\end{description}

\begin{thebibliography}{MVW}
\bibitem[GL]{GL} L.~Sernau, \emph{GRH for Dirichlet \(L\)-functions, in Lean~4}, a shorter rewriting of the Lean proof, text by Claude, \url{https://github.com/lsernau/grh-lean}, commit \texttt{72fead1}, readable edition \texttt{grh/}.
\bibitem[MVW]{MVW} C.~M\oe glin, M.-F.~Vign\'eras and J.-L.~Waldspurger, \emph{Correspondances de Howe sur un corps \(p\)-adique}, Lecture Notes in Math.\ 1291, Springer, 1987.
\bibitem[We]{We} A.~Weil, \emph{Sur certains groupes d'op\'erateurs unitaires}, Acta Math.\ \textbf{111} (1964), 143--211.
\end{thebibliography}

\end{document}
